% GENERATED FILE. Edit canonical lessons, then run npm run generate:books.
% canonical-source-hash: fnv1a64:ef51687a61bf57ef
% canonical-lessons: SA-C355-samsad, SA-C355-sadasyah, SA-C355-pratinidhih, SA-C355-adhyaksah, SA-C355-mantrayate, SA-R355-review-348-355

\chapter{The Assembly}
\label{ch:sa-the-assembly}
\hlchaptermodality{355}

\begin{chapteropening}
\textbf{By the end of this chapter:} \emph{I can say who sits in a parliament and what they do there.}

Complete the last lesson of chapter 355: I can say who sits in a parliament and what they do there.
\end{chapteropening}

\section[\texorpdfstring{\sk{संसद्} (saṃsad)}{संसद् (saṃsad)}]{\sk{संसद्} (saṃsad) — a parliament; an assembly}

\label{lesson:SA-C355-samsad}



\begin{quote}
Before the new one: say the Sanskrit for independent, then the Sanskrit for organizes.
\end{quote}

\subsection*{You'll want to know: \sk{संसद्}}
\textbf{\sk{संसद्}} — \emph{saṃsad} — \textquotedblleft{}a parliament; an assembly\textquotedblright{}.

\begin{quote}
\sk{संसत्} \sk{नूतनं} \sk{नियमं} \sk{करोति।}
\end{quote}

\begin{quote}
\emph{saṃsat nūtanaṃ niyamaṃ karoti.}
\end{quote}

\begin{quote}
Parliament makes a new law.
\end{quote}

A feminine noun.

The stem ends in \emph{-d}; said on its own, as here, the word ends in \emph{-t}: \sk{संसत्}, \emph{saṃsat}.

\begin{cousinweb}
\textbf{\sk{संसद्}} (\emph{saṃsad}) is built on \emph{sam}, \textquotedblleft{}together\textquotedblright{}, and the root √sad, \textquotedblleft{}to sit\textquotedblright{}; MW gives \emph{saṃ-sad} \textquotedblleft{}an assembly, a meeting, a court\textquotedblright{}, and \textquotedblleft{}a parliament\textquotedblright{} is its modern sense (MW).
\end{cousinweb}

\begin{grammarlens}[title={What you've built}]
One more word for giving reasons and opinions.
\end{grammarlens}

\subsection*{Your turn}
\begin{itemize}
\raggedright
  \item \emph{Say it:} \emph{saṃsad}
  \item \emph{Say it:} \emph{saṃsad}, once more
  \item \emph{Say it:} \emph{saṃsat nūtanaṃ niyamaṃ karoti.} — Parliament makes a new law.
  \item \emph{Recall:} say the Sanskrit for independent, then the Sanskrit for organizes, then say \emph{saṃsad} again
\end{itemize}

\begin{tcolorbox}[breakable,colback=teal!4,colframe=teal!35!black,title={Before you move on}]
What does \sk{संसद्} mean? (\textquotedblleft{}A parliament; an assembly\textquotedblright{}.) Say it once more.
\end{tcolorbox}
\section[\texorpdfstring{\sk{सदस्यः} (sadasyaḥ)}{सदस्यः (sadasyaḥ)}]{\sk{सदस्यः} (sadasyaḥ) — a member}

\label{lesson:SA-C355-sadasyah}



\begin{quote}
Before the new one: say the Sanskrit for organizes, then the Sanskrit for a parliament.
\end{quote}

\subsection*{You'll want to know: \sk{सदस्यः}}
\textbf{\sk{सदस्यः}} — \emph{sadasyaḥ} — \textquotedblleft{}a member\textquotedblright{}.

\begin{quote}
\sk{अरुणः} \sk{सभायाः} \sk{सदस्यः} \sk{अस्ति।}
\end{quote}

\begin{quote}
\emph{aruṇaḥ sabhāyāḥ sadasyaḥ asti.}
\end{quote}

\begin{quote}
Arun is a member of the assembly.
\end{quote}

A masculine noun.

\begin{cousinweb}
\textbf{\sk{सदस्यः}} (\emph{sadasyaḥ}) is built on \emph{sadas}, \textquotedblleft{}a seat, an assembly\textquotedblright{}, from the root √sad, \textquotedblleft{}to sit\textquotedblright{}; MW gives \emph{sadasya} \textquotedblleft{}an assessor\textquotedblright{} and \textquotedblleft{}a member of an assembly\textquotedblright{} (MW).
\end{cousinweb}

\begin{grammarlens}[title={What you've built}]
One more word for giving reasons and opinions.
\end{grammarlens}

\subsection*{Your turn}
\begin{itemize}
\raggedright
  \item \emph{Say it:} \emph{sadasyaḥ}
  \item \emph{Say it:} \emph{sadasyaḥ}, once more
  \item \emph{Say it:} \emph{aruṇaḥ sabhāyāḥ sadasyaḥ asti.} — Arun is a member of the assembly.
  \item \emph{Recall:} say the Sanskrit for organizes, then the Sanskrit for a parliament, then say \emph{sadasyaḥ} again
\end{itemize}

\begin{tcolorbox}[breakable,colback=teal!4,colframe=teal!35!black,title={Before you move on}]
What does \sk{सदस्यः} mean? (\textquotedblleft{}A member\textquotedblright{}.) Say it once more.
\end{tcolorbox}
\section[\texorpdfstring{\sk{प्रतिनिधिः} (pratinidhiḥ)}{प्रतिनिधिः (pratinidhiḥ)}]{\sk{प्रतिनिधिः} (pratinidhiḥ) — a representative}

\label{lesson:SA-C355-pratinidhih}



\begin{quote}
Before the new one: say the Sanskrit for a parliament, then the Sanskrit for a member.
\end{quote}

\subsection*{You'll want to know: \sk{प्रतिनिधिः}}
\textbf{\sk{प्रतिनिधिः}} — \emph{pratinidhiḥ} — \textquotedblleft{}a representative\textquotedblright{}.

\begin{quote}
\sk{प्रतिनिधिः} \sk{ग्रामस्य} \sk{मतं} \sk{वदति।}
\end{quote}

\begin{quote}
\emph{pratinidhiḥ grāmasya mataṃ vadati.}
\end{quote}

\begin{quote}
The representative speaks for the village.
\end{quote}

A masculine noun.

\begin{cousinweb}
\textbf{\sk{प्रतिनिधिः}} (\emph{pratinidhiḥ}) is built on \emph{prati}, \emph{ni} and the root √dhā, \textquotedblleft{}to put\textquotedblright{}; MW gives \emph{prati-nidhi} \textquotedblleft{}a substitute, a representative, a proxy\textquotedblright{} (MW).
\end{cousinweb}

\begin{grammarlens}[title={What you've built}]
One more word for giving reasons and opinions.
\end{grammarlens}

\subsection*{Your turn}
\begin{itemize}
\raggedright
  \item \emph{Say it:} \emph{pratinidhiḥ}
  \item \emph{Say it:} \emph{pratinidhiḥ}, once more
  \item \emph{Say it:} \emph{pratinidhiḥ grāmasya mataṃ vadati.} — The representative speaks for the village.
  \item \emph{Recall:} say the Sanskrit for a parliament, then the Sanskrit for a member, then say \emph{pratinidhiḥ} again
\end{itemize}

\begin{tcolorbox}[breakable,colback=teal!4,colframe=teal!35!black,title={Before you move on}]
What does \sk{प्रतिनिधिः} mean? (\textquotedblleft{}A representative\textquotedblright{}.) Say it once more.
\end{tcolorbox}
\section[\texorpdfstring{\sk{अध्यक्षः} (adhyakṣaḥ)}{अध्यक्षः (adhyakṣaḥ)}]{\sk{अध्यक्षः} (adhyakṣaḥ) — a president; a chairperson}

\label{lesson:SA-C355-adhyaksah}



\begin{quote}
Before the new one: say the Sanskrit for a member, then the Sanskrit for a representative.
\end{quote}

\subsection*{You'll want to know: \sk{अध्यक्षः}}
\textbf{\sk{अध्यक्षः}} — \emph{adhyakṣaḥ} — \textquotedblleft{}a president; a chairperson\textquotedblright{}.

\begin{quote}
\sk{अध्यक्षः} \sk{सभाम्} \sk{आरभते।}
\end{quote}

\begin{quote}
\emph{adhyakṣaḥ sabhām ārabhate.}
\end{quote}

\begin{quote}
The president opens the meeting.
\end{quote}

A masculine noun.

\begin{cousinweb}
\textbf{\sk{अध्यक्षः}} (\emph{adhyakṣaḥ}) is built from \emph{adhi}, \textquotedblleft{}over\textquotedblright{}, and \emph{akṣa}, \textquotedblleft{}the eye\textquotedblright{}; MW gives \emph{adhy-akṣa} \textquotedblleft{}an inspector, a superintendent\textquotedblright{}, and \textquotedblleft{}a president\textquotedblright{} is its modern sense (MW).
\end{cousinweb}

\begin{grammarlens}[title={What you've built}]
One more word for giving reasons and opinions.
\end{grammarlens}

\subsection*{Your turn}
\begin{itemize}
\raggedright
  \item \emph{Say it:} \emph{adhyakṣaḥ}
  \item \emph{Say it:} \emph{adhyakṣaḥ}, once more
  \item \emph{Say it:} \emph{adhyakṣaḥ sabhām ārabhate.} — The president opens the meeting.
  \item \emph{Recall:} say the Sanskrit for a member, then the Sanskrit for a representative, then say \emph{adhyakṣaḥ} again
\end{itemize}

\begin{tcolorbox}[breakable,colback=teal!4,colframe=teal!35!black,title={Before you move on}]
What does \sk{अध्यक्षः} mean? (\textquotedblleft{}A president; a chairperson\textquotedblright{}.) Say it once more.
\end{tcolorbox}
\section[\texorpdfstring{\sk{मन्त्रयते} (mantrayate)}{मन्त्रयते (mantrayate)}]{\sk{मन्त्रयते} (mantrayate) — consults; deliberates}

\label{lesson:SA-C355-mantrayate}



\begin{quote}
Before the new one: say the Sanskrit for a representative, then the Sanskrit for a president.
\end{quote}

\subsection*{You'll want to know: \sk{मन्त्रयते}}
\textbf{\sk{मन्त्रयते}} — \emph{mantrayate} — \textquotedblleft{}consults; deliberates\textquotedblright{}.

\begin{quote}
\sk{मीरा} \sk{मित्रेण} \sk{सह} \sk{मन्त्रयते।}
\end{quote}

\begin{quote}
\emph{mīrā mitreṇa saha mantrayate.}
\end{quote}

\begin{quote}
Mira talks it over with her friend.
\end{quote}

Past (he …): \sk{सः} \sk{मित्रेण} \sk{सह} \sk{मन्त्रितवान्}, \emph{saḥ mitreṇa saha mantritavān}.

\begin{cousinweb}
\textbf{\sk{मन्त्रयते}} (\emph{mantrayate}) is built on \emph{mantra}, \textquotedblleft{}a sacred text, counsel\textquotedblright{}; MW gives \emph{mantrita} \textquotedblleft{}discussed, deliberated, determined\textquotedblright{} (MW).
\end{cousinweb}

\begin{grammarlens}[title={What you've built}]
One more word for giving reasons and opinions.
\end{grammarlens}

\subsection*{Your turn}
\begin{itemize}
\raggedright
  \item \emph{Say it:} \emph{mantrayate}
  \item \emph{Say it:} \emph{mantrayate}, once more
  \item \emph{Say it:} \emph{mīrā mitreṇa saha mantrayate.} — Mira talks it over with her friend.
  \item \emph{Recall:} say the Sanskrit for a representative, then the Sanskrit for a president, then say \emph{mantrayate} again
\end{itemize}

\begin{tcolorbox}[breakable,colback=teal!4,colframe=teal!35!black,title={Before you move on}]
What does \sk{मन्त्रयते} mean? (\textquotedblleft{}Consults; deliberates\textquotedblright{}.) Say it once more.
\end{tcolorbox}
\section[(dialogue)]{Looking back: from Fair and Unfair to The Assembly}

\label{lesson:SA-R355-review-348-355}



\begin{quote}
Say the Sanskrit for exploitation and fair.
\end{quote}

\subsection*{Your turn}
Say these aloud:

\begin{itemize}
\raggedright
  \item exploitation, fair, a precept, partiality, breaks (a rule)
  \item a court of law, a defendant, does wrong, restrains, an oath
  \item captivity, release, settles (a question), gives back, atonement
  \item a virtuous act, discernment, conforms to, abides by, the bounds of propriety
  \item violence, tolerant, does harm to, shows favour to, forbearance
  \item government, presides over, a nation, a governor, an official order
  \item a citizen (of a country), the public, society, independent, organizes
  \item a parliament, a member, a representative, a president, consults
\end{itemize}

\begin{tcolorbox}[breakable,colback=teal!4,colframe=teal!35!black,title={Before you move on}]
Exploitation? (\textbf{\sk{शोषणम्}}, \emph{śoṣaṇam}.) Fair? (\textbf{\sk{न्याय्यः}}, \emph{nyāyyaḥ}.) A precept? (\textbf{\sk{विधिः}}, \emph{vidhiḥ}.) Partiality? (\textbf{\sk{पक्षपातः}}, \emph{pakṣapātaḥ}.) Breaks (a rule)? (\textbf{\sk{उल्लंघयति}}, \emph{ullaṅghayati}.) A court of law? (\textbf{\sk{न्यायालयः}}, \emph{nyāyālayaḥ}.) A defendant? (\textbf{\sk{प्रतिवादी}}, \emph{prativādī}.) Does wrong? (\textbf{\sk{अपराध्यति}}, \emph{aparādhyati}.) Restrains? (\textbf{\sk{निगृह्णाति}}, \emph{nigṛhṇāti}.) An oath? (\textbf{\sk{शपथः}}, \emph{śapathaḥ}.) Captivity? (\textbf{\sk{बन्धनम्}}, \emph{bandhanam}.) Release? (\textbf{\sk{मुक्तिः}}, \emph{muktiḥ}.) Settles (a question)? (\textbf{\sk{निर्धारयति}}, \emph{nirdhārayati}.) Gives back? (\textbf{\sk{प्रत्यर्पयति}}, \emph{pratyarpayati}.) Atonement? (\textbf{\sk{प्रायश्चित्तम्}}, \emph{prāyaścittam}.) A virtuous act? (\textbf{\sk{सत्कर्म}}, \emph{satkarma}.) Discernment? (\textbf{\sk{विवेकः}}, \emph{vivekaḥ}.) Conforms to? (\textbf{\sk{अनुवर्तते}}, \emph{anuvartate}.) Abides by? (\textbf{\sk{अनुसरति}}, \emph{anusarati}.) The bounds of propriety? (\textbf{\sk{मर्यादा}}, \emph{maryādā}.) Violence? (\textbf{\sk{हिंसा}}, \emph{hiṃsā}.) Tolerant? (\textbf{\sk{सहिष्णुः}}, \emph{sahiṣṇuḥ}.) Does harm to? (\textbf{\sk{अपकरोति}}, \emph{apakaroti}.) Shows favour to? (\textbf{\sk{अनुगृह्णाति}}, \emph{anugṛhṇāti}.) Forbearance? (\textbf{\sk{क्षमा}}, \emph{kṣamā}.) Government? (\textbf{\sk{शासनम्}}, \emph{śāsanam}.) Presides over? (\textbf{\sk{अधितिष्ठति}}, \emph{adhitiṣṭhati}.) A nation? (\textbf{\sk{राष्ट्रम्}}, \emph{rāṣṭram}.) A governor? (\textbf{\sk{शासकः}}, \emph{śāsakaḥ}.) An official order? (\textbf{\sk{आदेशः}}, \emph{ādeśaḥ}.) A citizen (of a country)? (\textbf{\sk{नागरिकः}}, \emph{nāgarikaḥ}.) The public? (\textbf{\sk{जनता}}, \emph{janatā}.) Society? (\textbf{\sk{समाजः}}, \emph{samājaḥ}.) Independent? (\textbf{\sk{स्वतन्त्रः}}, \emph{svatantraḥ}.) Organizes? (\textbf{\sk{संघटयति}}, \emph{saṅghaṭayati}.) A parliament? (\textbf{\sk{संसद्}}, \emph{saṃsad}.) A member? (\textbf{\sk{सदस्यः}}, \emph{sadasyaḥ}.) A representative? (\textbf{\sk{प्रतिनिधिः}}, \emph{pratinidhiḥ}.) A president? (\textbf{\sk{अध्यक्षः}}, \emph{adhyakṣaḥ}.) Consults? (\textbf{\sk{मन्त्रयते}}, \emph{mantrayate}.)
\end{tcolorbox}
